%% kapa: exam=1 round=24 year=2013 period=2 form=A subject=appraisal_law subject_ko=부동산관계법규 seq=23 number=63
%% source: Q-Net 감정평가사 기출문제 (cst00309) artlSeq=5073802 file=extracted/1008/5073802/2013년 제24회 감정평가사 1차시험 기출문제/2교시(A형).hwp
%% conversion: kordoc 4.24.51 -> tools/gen1.py (auto); review=proofread-sonnet; flags=none
\begin{question}{63}
부동산 가격공시 및 감정평가에 관한 법령상 표준지공시지가에 대한 이의신청에 관한 설명으로 옳지 않은 것을 모두 고른 것은?
\begin{framedtext}
ㄱ. 표준지공시지가에 대한 이의신청은 구두로 할 수 있다. \par
ㄴ. 표준지공시지가에 대한 이의신청은 국토교통부장관에게 한다. \par
ㄷ. 이의신청에 대한 심사는 이의신청이 있은 날부터 30일 이내에 하여야 한다. \par
ㄹ. 토지소유자가 아닌 자는 이의신청을 할 수 없다. \par
ㅁ. 중앙부동산평가위원회는 표준지공시지가에 대한 이의신청에 관한 사항을 심의한다.
\end{framedtext}
\begin{choices}
  \item ㄱ, ㄴ, ㄷ
  \item ㄱ, ㄷ, ㄹ
  \item ㄱ, ㄹ, ㅁ
  \item ㄴ, ㄷ, ㅁ
  \item ㄴ, ㄹ, ㅁ
\end{choices}
\end{question}
\answer{2}
